% ID: SEL-D0005-OB-E04
% ESTADO: cotejada. Segunda pasada visual completa y borrador compilado revisados.
% Fuente: Selectividad/sel_2019_matematicas_aplicadas/Reserva A/Exámenes Mat Aplicadas CCSS - A y B.pdf
% SHA-256: d28871f7f203cf1be9e479e9d642cdb6b1623704a8b76a2b57cfa2d12950f293
% Página: 2; opción B; ejercicio 4.

El tiempo de duración, en horas, de un modelo de bombilla LED, sigue una ley
Normal de media desconocida y desviación típica 150 horas. Con una muestra de
bombillas de ese modelo y a un nivel de confianza del 98.5\,\% se ha obtenido
que el intervalo de confianza para la media es $(18\,475.7,\ 18\,524.3)$.
\begin{enumerate}[label=\alph*)]
  \item \textbf{(1.5 puntos)} Calcule el valor que se obtuvo para la media
  de la muestra y el tamaño de la muestra utilizado.
  \item \textbf{(1 punto)} \textquestiondown{}Cuál será el error máximo de
  estimación de la media si se hubiese utilizado una muestra de tamaño 100
  y un nivel de confianza del 96.6\,\%?
\end{enumerate}
